% !TeX root = ../../Book2.tex
\section{常系数算子的基本解}
\label{b2:sec:0504}

常系数微分在频率侧变成多项式乘法。若能解出对点质量的响应，就可以通过卷积生成其他右端的解。这个响应称为基本解；它通常带有奇性，因此分布语言不可省略。

\begin{definition}\label{b2:def:05-fundamental}
设 \(P(\partial)=\sum_{|\alpha|\leq m}a_\alpha\partial^\alpha\) 为常系数微分算子。若分布 \(E\) 满足 \(P(\partial)E=\delta_0\)，则称 \(E\) 为该算子的\term[jiben-jie]{基本解}[fundamental solution]。
\end{definition}

\subsection{微分算子的符号}
\label{b2:subsec:050401}

定义其\term[weifen-suanzi-fuhao]{符号}[symbol]为
\[
 p(\xi)\coloneqq\sum_{|\alpha|\leq m}a_\alpha(2\pi\mathrm{i}\xi)^\alpha.
\]
由\eqref{b2:eq:05-distribution-operations}，\(\mathcal F(P(\partial)T)=p\widehat T\)。例如 \(-\Delta\) 的符号为 \(4\pi^2|\xi|^2\)，而 \(1-\Delta\) 的符号为 \(1+4\pi^2|\xi|^2\)。正号不同会直接改变能量与基本解，不能只记住“Laplacian 对应频率平方”。

\subsection{Fourier 方法}
\label{b2:subsec:050402}

若基本解温和，则其频率方程为 \(p\widehat E=1\)。当 \(p\) 处处非零，且 \(1/p\) 及其全部导数至多多项式增长时，可以定义 \(\widehat E=1/p\)，再取分布逆变换。这给出一个充分条件。若 \(p\) 有零点，就不能直接在这些点除法；即使几乎处处可以相除，也须检查所得函数是否局部可积。

\begin{proposition}\label{b2:prop:05-bessel-fundamental}
令
\[
 H_t(x)=(4\pi t)^{-d/2}e^{-|x|^2/(4t)},\quad
 E(x)=\int_0^\infty e^{-t}H_t(x)\,\mathrm{d}t.
\]
则 \(E\in L^1(\mathbb R^d)\)，非负且积分为 \(1\)，并在分布意义下满足 \((1-\Delta)E=\delta_0\)。
\end{proposition}
\begin{proof}
Gauss 积分给出 \(\int H_t=1\)，所以 Tonelli 说明 \(\int E=\int_0^\infty e^{-t}\,\mathrm{d}t=1\)，并保证定义积分几乎处处有限。由缩放公式，\(\widehat H_t(\xi)=e^{-4\pi^2t|\xi|^2}\)。再次按绝对积分 \(\int e^{-t}\|H_t\|_1\,\mathrm{d}t=1\) 换序，得到
\[
 \widehat E(\xi)=\int_0^\infty e^{-t(1+4\pi^2|\xi|^2)}\,\mathrm{d}t
 =\dfrac1{1+4\pi^2|\xi|^2}.
\]
因此 \(\mathcal F((1-\Delta)E)=1=\widehat\delta_0\)，分布变换的单射性给出结论。
\end{proof}

在一维还可直接写出 \(E(x)=\tfrac12e^{-|x|}\)：它在原点以外满足 \(E''=E\)，而一阶导数在原点的跳跃为 \(-1\)，所以 \(E''=E-\delta_0\)。这与频率分母 \(1+4\pi^2\xi^2\) 一致，也说明基本解方程的点质量来自导数的跳跃。

\subsection{Laplace 算子的基本解}
\label{b2:subsec:050403}

以下统一求 \(-\Delta\) 的基本解。记 \(\sigma_{d-1}\) 为单位球面的面积；这里的符号是表面积，不是单位球体积。

\begin{theorem}\label{b2:thm:05-laplace-fundamental}
局部可积函数
\[
 \Phi_d(x)=
 \begin{cases}
 \dfrac1{(d-2)\sigma_{d-1}}|x|^{2-d},&d\geq3,\\[4pt]
 -\dfrac1{2\pi}\log|x|,&d=2,\\[4pt]
 -\dfrac12|x|,&d=1
 \end{cases}
\]
定义温和分布，满足 \(-\Delta\Phi_d=\delta_0\)。
\end{theorem}
\begin{proof}
原点附近，\(d\geq3\) 时的绝对积分由 \(\int_0^1r^{2-d}r^{d-1}\,\mathrm{d}r=1/2\) 控制；二维由 \(\int_0^1r|\log r|\,\mathrm{d}r<\infty\) 控制。一维函数连续。无穷远只有幂或对数增长，故均满足温和性的加权积分条件。

当 \(d\geq2\) 时，在原点外用径向公式 \(\Delta v=v''+(d-1)v'/r\) 验证 \(\Delta\Phi_d=0\)。取 \(\varphi\in C_c^\infty\)，在去掉半径 \(\varepsilon\) 小球的区域上使用 Green 公式。外边界处测试函数及导数为零；内边界的区域外法向是 \(-e_r\)，故
\[
 \int_{|x|>\varepsilon}\Phi_d(-\Delta\varphi)\,\mathrm{d}x
 =\int_{|x|=\varepsilon}
 \left(-\varphi\,\partial_r\Phi_d+\Phi_d\partial_r\varphi\right)\,\mathrm{d}S.
\]
在两种维数情形中，都有 \(-\partial_r\Phi_d(\varepsilon)=1/(\sigma_{d-1}\varepsilon^{d-1})\)。第一项正好是 \(\varphi\) 的球面平均，趋于 \(\varphi(0)\)。第二项的绝对值至多为常数乘 \(\varepsilon\)（\(d\geq3\)）或 \(\varepsilon|\log\varepsilon|\)（\(d=2\)），故趋零。左侧由局部可积性趋于整个积分，得到分布恒等式。

一维时，分别在负、正半轴作两次分部积分，有 \(\int |x|\varphi''(x)\,\mathrm{d}x=2\varphi(0)\)，因此 \(-\Phi_1''=\delta_0\)。首先得到的等式在 \(\mathcal D\) 上成立；两边均温和，利用 \(\mathcal D\) 在 \(\mathcal S\) 中稠密便推广到 \(\mathcal S'\)。
\end{proof}

二维的对数和一维的线性增长，正好提醒我们：\(1/(4\pi^2|\xi|^2)\) 在这些维数的原点并不可积，不能用该函数的普通积分冒充基本解的频率分布。上面的空间构造已经定义了合法的温和基本解，其变换满足 \(4\pi^2|\xi|^2\widehat\Phi_d=1\)；有奇性的频率乘法方程应在这个意义下理解。

\subsection{基本解与方程求解}
\label{b2:subsec:050404}

\begin{proposition}\label{b2:prop:05-fundamental-convolution}
设温和分布 \(E\) 满足 \(P(\partial)E=\delta_0\)，且 \(f\in C_c^\infty(\mathbb R^d)\)。则
\[
 u(x)\coloneqq\langle E_y,f(x-y)\rangle
\]
是光滑函数，并满足 \(P(\partial)u=f\)。
\end{proposition}
\begin{proof}
对固定 \(x\)，函数 \(y\mapsto f(x-y)\) 属于 \(\mathcal S\)。当 \(x\) 限于紧集时，这族函数有共同紧支集，并且关于 \(x\) 的差商在每个 Schwartz 半范数中趋于相应导数；这可对各阶导数使用 Taylor 公式的一致余项估计。连续泛函 \(E\) 可以通过这些极限，所以 \(u\) 光滑，导数可移到 \(f\) 上。

又因 \(\partial_x^\alpha f(x-y)=(-1)^{|\alpha|}\partial_y^\alpha f(x-y)\)，分布导数定义给出
\[
 P(\partial)u(x)=\langle P(\partial)E_y,f(x-y)\rangle
 =\langle\delta_0,f(x-\cdot)\rangle=f(x).
\]
\end{proof}

特别地，\(u=\Phi_d*f\) 解 \(-\Delta u=f\)。这只解决全空间的方程，没有自动满足某个有界区域上的 Dirichlet 或 Neumann 条件。基本解也一般不唯一：加上任何齐次方程的温和解仍为基本解。例如给 \(\Phi_d\) 加一个调和多项式不改变 \(-\Delta\Phi_d=\delta_0\)。要获得唯一的边值解，还须结合区域、边界条件与增长约束。



\ProseParagraphBreak

二维基本解还表现出临界尺度的差别。对 \(\lambda>0\)，
\[
 \Phi_2(\lambda x)=\Phi_2(x)-\dfrac{\log\lambda}{2\pi}.
\]
它不严格齐次，而是相差一个常数；该常数的 Laplace 为零，所以仍与基本解方程相容。高维的 \(\Phi_d\) 则按次数 \(2-d\) 齐次。这说明从一般维数公式形式代入 \(d=2\) 不能直接得到答案：分母与幂次同时退化，要保留对数结构及可加的齐次解。
